\section{Mathématiques : mode d'emploi}
\subsection{But}

Etablir de manière logique des orioruérés sur des objers en partant de 
\begin{itemize}
    \item propriétés déjà étables
    \item propriétés admises = axiomes
\end{itemize}

\subsection{Support matériel}

Les objets mathématiques sont représentés par des lettres, parfois déclinées avec des 
machins bizarres ou indexées. 

Règle fondamentale : deux objets différents portent deux noms différents. 

\subsection{Objets étudiés}

En maths actuelles, on étudie les ensembles. 

Relation fondamentale entre ensembles : l'appartenance. Si $a$ et $b$ sont des 
ensembles, deux cas possibles : 

\begin{itemize}
    \item $a$ appartient à $b$, noté $a \in b$
    \item $a$ n'appartient pas à $b$, noté $a \notin b$
\end{itemize}

\begin{definition}{Inclusion}{}
    Si $a$, $b$ sont des ensembles, on dit que $a$ est inclus dans $b$ et on note 
    $a \subset b$ ssi pour tout élément $x$ de $a$, on a $x \in b$
\end{definition}

\subsection{La théorie des ensembles}

\underline{Problème} : il faut choisir une axiomatique assez riche pour être utile, 
pas trop pour ne pas être contradictoire. 

\underline{Choix :} Théorème des ensembles de Zemelo Frankil avec l'axiome du choix.

En gros : 

\begin{enumerate}
    \item Deux ensembles sont identiques si ils ont les même éléments. 
    \item Il existe un ensemble, l'ensemble vide, ne contenant aucun élément
    \item Si $a$ et $b$ sont des ensembles, il existe un ensemble $c = \{a, b\}$
            dont les éléments sont exactement $a$ et $b$
    \item Réunion : si $a$ est un ensemble, il existe un ensemble $b$ dont les éléments
        sont les éléments des éléments de $a$. Donc $b = \bigcup_{x \in a} x$.
    \item Ensemble des parties : si $a$ est un ensemble, il existe un ensemble noté
        $\mathbb P(a)$ dont les éléments sont les sous ensembles de $a$
    \item Produit cartésien : Si $a$ et $b$ sont deux ensembles, il existe un ensemble
        noté $a \times b$ dont les éléments sont les couples $(x, y)$ où $x \in a$ 
        et $y \in b$
    \item Axiome de compréhension : Si $a$ est un ensemble et $P$ une propriété réalisable
        par des éléments de $a$, on peut constituer l'ensemble $b = \{x \in a \mid P(x)\}$
    \item Axiome de l'infini : L'ensemble $\mathbb N$ existe.
    \item Axiome du choix : Si $A$ est un ensemble d'ensembles non vides, il existe une 
        application $\fonction{f}{A}{E}{P}{x}$ avec $x \in P$ où $E = \bigcup_{P \in A} P$
\end{enumerate}

\section{Les types de raisonements}

\subsection{Règle générale}

C'est le réusltat à démontrer qui guide dans le choix du raisonnement utilisé et des
acteurs mis en oeuvre. 

\subsection{Types d'énoncés}
\subsubsection{Cas quelconques}

Raisonnement par l'absurde (ne pas utiliser en priorité)

\subsubsection{Enoncés du type P ou Q }
On suppose que $P$ n'est pas réalisé et on montre $Q$.

\subsubsection{Enoncé P implique Q }

On suppose $P$, on en déduit par une suite de donc que $Q$ l'est aussi.

On peut également raisonner par contraposée, supposer non Q et en déduire non P. 

\subsubsection{Enoncé P equivalent à Q }

On montre les deux implications.

\subsubsection{Enoncé A = B }

On montre les deux inclusions. 

\subsubsection{Enoncés du type pour tout x de E, P(x)}

On considère un élément $x$ de $E$ et on établit $P(x)$. 

\subsubsection{Enoncé du type il existe un unique x dans E, P(x)}

\underline{Analyse - synthèse}:
\begin{itemize}
    \item Analyse : On suppose qu'un tel $x$ existe et on l'identifie. On en 
        déduit l'uniciré en cas d'existence. 
    \item Synthèse : on vérifie que la valeur identifiée convient. 
\end{itemize}


\underline{Existence - unicité}
\begin{itemize}
    \item Exitence : on exhibe un $x$ qui convient
    \item unicité : on suppose que $x_1$ et $x_2$ conviennent, et on montre que 
        $x_1 = x_2$
\end{itemize}

\subsubsection{Enoncé du type il existe x dans E, P(x)}

On devine et exhibe un $x$ qui convient ou on commence une analyse jusqu'à obtenir
suffisament d'informations pour identifier un tel $x$. 

\subsubsection{Enoncé pour tout n dans N, P(n)}

On peut raisonner par récurrence : 
\begin{itemize}
    \item simple
    \item double
    \item avec prédécesseurs (= forte)
\end{itemize}

